\clearpage
\begin{summarybox}
\textbf{C085\quad 双曲线焦点平行渐近线截线定比——快记／快懂／快用}

\textbf{快记（一句话＋两条可迁移关系）：}
\begin{center}
\textbf{“焦点过线平行渐近线；近段比远段，等于斜率平方。”}
\end{center}
这里“近段”专指 $FM$，“远段”专指 $MN$，点的顺序为 $F-M-N$。
\[
\boxed{\lambda=\frac{FM}{MN}
=\left(\frac ba\right)^2=e^2-1},
\qquad\boxed{e=\sqrt{1+\lambda}}.
\]
\[
\boxed{\frac{FN}{MN}=e^2},\qquad
\boxed{\frac{FM}{FN}=1-\frac1{e^2}}.
\]

\textbf{快懂（最低门槛三连）：}
\[
\begin{gathered}
F=(c,0),\quad N_x=\frac c2,\quad
M=\left(c-t,-\frac ba t\right),\\[2pt]
(c-t)^2-t^2=a^2
\Longrightarrow t=\frac{b^2}{2c},\\[2pt]
\frac{FM}{MN}=\frac{t}{c/2-t}=\frac{b^2}{a^2}.
\end{gathered}
\]
\textbf{为什么出现平方差：}直线斜率恰好等于渐近线斜率，
代入后 $y^2/b^2$ 中的 $b^2$ 被抵消。

\textbf{快用（最少步骤）：}
\begin{itemize}
\item 已知 $FM:MN=\lambda:1$：直接 $e=\sqrt{1+\lambda}$；
\item 已知 $MN:FM=\mu:1$：先取倒数，再 $e=\sqrt{1+1/\mu}$；
\item 已知 $FN:MN=\nu:1$：直接 $e=\sqrt{\nu}$；
\item 已知渐近线的斜率 $k$：$\lambda=k^2$、$e=\sqrt{1+k^2}$；
\item 要求线段\textbf{绝对长度}：$FM=b^2/(2a)$、$MN=a/2$、$FN=c^2/(2a)$。
\end{itemize}

\textbf{证明方法如何选：}
\textbf{最快}——因式分解＋线性分点（不求交点）；
\textbf{门槛最低}——只设一个 $t$＋平方差；
\textbf{常规}——列直线、求 $M,N$、算距离；
\textbf{拓展}——斜坐标 $UV=1$，三个点的 $V$ 为 $e,1/e,0$。

\textbf{一眼防错：}
\boxed{F-M-N}\quad\text{、}\quad
\boxed{\lambda=e^2-1}\quad\text{、}\quad
\boxed{c^2=a^2+b^2}.

\textbf{思维迁移：}看到“\emph{平行圆锥曲线的渐近线}”，
优先尝试\emph{平方差因式分解或把一个二次项消掉}，
往往能把交点问题降为一次问题。
\end{summarybox}
